% Suggested final budget: 1.0 page, including the task table.
\section{Tasks and Expert Demonstrations}
\label{sec:data}
\subsection{Task coverage}
% Sources: dp-manip/docs/final-plan.md section 1 (planned tasks and the
% pre-registered LiftPegUpright fallback), dp-manip/docs/experiment-progress.zh-CN.md
% section 5.1 (PlaceSphere rationale), maniskill-demogen/README.md (versions,
% observation layout).
We use six ManiSkill~3.0.1 tabletop tasks for the Panda arm (\cref{tab:tasks})
that cover
prehensile and non-prehensile manipulation: lifting, stacking, pushing,
pulling, placing into a bin, and reorienting a peg. PushCube and PullCube
both move a cube without grasping it but need opposite contact strategies:
pushing from the near side versus reaching behind the cube and dragging it
back. Policies see
$128\times128$ RGB images and the low-dimensional state that ManiSkill
exposes in RGB observation mode (robot state and task extras such as the goal
position); object poses are never given to the policy. Every task uses one
fixed camera, and StackCube adds a wrist camera.

\textbf{Changes from the proposal.} The proposal listed PegInsertionSide and
PlugCharger as its two precision tasks. We could not establish a reliable
baseline on either, so both were replaced and neither appears in the results
below. LiftPegUpright, the fallback named in the proposal, replaces
PlugCharger. PlaceSphere replaces PegInsertionSide; it was chosen because the
failure-correction study needs a baseline with both successful and failed
rollouts.
% Sources (all to be confirmed against the evaluated runs):
% - control mode, horizon: dp-manip/configs/tasks/*.toml. LiftPegUpright's
%   horizon is left open: the config default (50) differs from the horizons
%   used in evaluation (200 and 400).
% - action dim: dp-manip/docs/results-summary.zh-CN.md section 1.
% - 400 / 50 demonstrations: the exporter keeps exactly the first 400 / 50
%   usable demonstrations per split (maniskill-demogen/tasks.py) and fails if
%   there are fewer; confirm the counts in each export_info.json.
% - demo length: PlaceSphere from dp-manip/docs/experiment-progress.zh-CN.md
%   section 6; LiftPegUpright from results-summary section 5.
% - replay kept: PlaceSphere 298 of 440 in the first generation
%   (task-owner-runbook section 2). PegInsertionSide's 97% is a pilot figure
%   (34 of 35, maniskill-demogen/tasks.py); replace it with the final log. The
%   cube tasks were close to 100% in pilots and are left open.
% - success rule: dp-manip/docs/final-plan.md section 1 for the original
%   tasks; PlaceSphere and LiftPegUpright from the ManiSkill task documentation
%   (maniskill.readthedocs.io, latest version, read 2026-10-12); confirm that
%   version 3.0.1 uses the same conditions.
\begin{table}[htbp]
  \centering
  \caption{Task coverage (ManiSkill \texttt{-v1} environments), with the
  action dimension in parentheses. Horizon is the evaluation episode limit and
  demo length the training-episode length, in steps. Replay kept is the share
  of expert trajectories that pass the export checks. Every task has 400
  training and 50 validation demonstrations.}
  \label{tab:tasks}
  \small
  \setlength{\tabcolsep}{3.5pt}
  \begin{tabularx}{\linewidth}{@{}llcccc>{\raggedright\arraybackslash}X@{}}
    \toprule
    Task & Skill & Control & Horizon &
    \shortstack{Demo length\\median (range)} & \shortstack{Replay\\kept} &
    Success rule \\
    \midrule
    PickCube & Grasp, lift & EE $\Delta$pos (4) & 100 & \pending & \pending & Cube within 2.5 cm of goal, robot static \\
    StackCube & Grasp, stack & EE $\Delta$pos (4) & 200 & \pending & \pending & Cube A on cube B, released, static \\
    PushCube & Push & EE $\Delta$pos (4) & 100 & \pending & \pending & Cube inside goal region \\
    PullCube & Pull & EE $\Delta$pos (4) & 100 & \pending & \pending & Cube inside goal region \\
    PlaceSphere & Grasp, place & EE $\Delta$pos (4) & 200 & 113 (90--150) & 68\% & Sphere on the bin, robot static, gripper open \\
    LiftPegUpright & Reorient & Joint pos (8) & \pending & 169 (153--702) & \pending & Peg upright within 0.08 rad, resting on the table \\
    \midrule
    \multicolumn{7}{@{}l}{\textit{Attempted, reported as a negative result; not one of the six}} \\
    PegInsertionSide & Insert & Joint pos (8) & 300 & \pending & 97\% & Peg head inside hole \\
    \bottomrule
  \end{tabularx}
\end{table}


\subsection{Demonstration generation and quality}
% Drafted 2026-10-12. Sources: maniskill-demogen/README.md, docs/pipeline.md,
% tasks.py and scripts/run_cpu.py (pipeline, checks, control modes);
% dp-manip/docs/experiment-progress.zh-CN.md section 6 and
% docs/task-owner-runbook.zh-CN.md section 2 (PlaceSphere yield); the
% LiftPegUpright member report task6_backbone, section 2.1 (its cameras and
% episode lengths); maniskill-demogen/tasks.py (440 / 55 default counts).
We generated every demonstration ourselves with the motion-planning solvers
that ManiSkill provides for the Panda arm; no dataset was downloaded. The
solver is run on consecutive environment seeds and only successful episodes
are kept. Each is then replayed twice in the policy's control mode on the same
CPU physics backend, once recording $128\times128$ RGB observations and once
recording state. The exporter accepts an episode only if both replays succeed
and agree step by step, drops episodes that contain a joint jump above
0.2~rad, and recomputes the first frame from a fresh reset, because
ManiSkill~3.0.1 records stale contact observations there. It writes the first
400 training and 50 validation episodes in seed order and fails if fewer are
usable.

The solver acts in joint-position control. The cube tasks and PlaceSphere are
converted to end-effector delta-position control. LiftPegUpright needs
end-effector rotation; in pilots on rotation tasks, conversion to delta-pose
control kept only 60--73\% of the trajectories, below the 90\% threshold we
had set, so it keeps joint-position control.

By default the generator produces 440 training and 55 validation expert
trajectories per task and exports the first 400 and 50 that pass the checks.
PlaceSphere was the exception: only 68\% of its expert trajectories passed
(298 of 440), so we regenerated it from 700. Episode length differs widely between tasks:
PlaceSphere demonstrations take 90--150 steps (median 113), while those of
LiftPegUpright take 153--702 (median 169, first 100 demonstrations).

\subsection{Data separation}
\label{sec:separation}
% Drafted 2026-10-12. Sources: maniskill-demogen/tasks.py and docs/pipeline.md
% (demonstration seed ranges, export assertions), dp-manip/configs/baseline.toml
% [eval] (validation and test seeds), dp-manip/STATUS.md phase 4 (nested
% subsets, duplicate check, recorded demo seeds), dp-manip/dp_manip/data.py
% (normalization), configs/failure_aware/takeover_*.toml (takeover seeds).
Environment seeds define four disjoint roles. Training demonstrations come
from seeds 0--3999 and validation demonstrations from 4000--4999; validation
rollouts use seeds from 5000, and the held-out test episodes use
10000--10099. Test seeds are never used for training, for the validation
loss, or to choose a checkpoint, since we always report the final one. The
takeover-correction study uses further ranges that overlap none of these:
40000--40949 for its probes, 41000--41599 to collect corrections, and
43000--43199 for its evaluation.

A data-size subset is the first $N$ training demonstrations in seed order.
Subsets are therefore nested and identical across training seeds, and the
action and state normalization statistics are computed from the selected $N$
demonstrations only. The exporter asserts the seed ranges, and the training
code rejects duplicate episodes and records the selected seeds with every run.
